% =============================================================================================
\chapter*{Résumé}
\addcontentsline{toc}{chapter}{Résumé}

Une banque qui finance des projets d'investissement de PME doit arbitrer entre des critères contradictoires
(risque, rentabilité, capacité de remboursement, garanties, impact environnemental et social) sous un budget
limité. Ce projet construit un système d'aide à la décision dans lequel chaque composant a un rôle unique :
\textbf{l'apprentissage automatique estime le risque, une méthode multicritère classe et catégorise, un agent
d'arbitrage respecte le budget, un modèle de langage rédige et répond, et le comité d'investissement décide.}

Le risque de défaut à cinq ans est appris sur 910\,033 prêts réels du programme américain SBA 7(a) (2000-2016),
séparés dans le temps (entraînement 2000-2012, validation 2013-2014, test 2015-2016). Une fuite de données a été
détectée et corrigée : la durée exacte du prêt est modifiée après un défaut. Trois modèles sont comparés sur les
mêmes 13 variables : une régression logistique (AUROC de 0,726 sur le test), un modèle additif interprétable EBM
(0,778) et un modèle de boosting LightGBM (0,800). La règle, fixée avant l'entraînement, retient le meilleur modèle
\emph{interprétable} : l'EBM, calibré par régression isotonique. L'interprétabilité coûte ainsi 0,022 point d'AUROC.

Quarante prêts du test, dont les indicateurs financiers sont simulés à partir des paramètres sectoriels de
Damodaran, forment le portefeuille candidat. Les poids des huit critères sont obtenus par AHP (ratio de cohérence
0,009) ; un filtre éliminatoire écarte 11 projets ; ELECTRE TRI affecte 9 projets à \og Financer \fg{}, 19 à
\og Financer sous conditions \fg{} et 12 à \og Refuser \fg{} ; ELECTRE III et PROMETHEE II donnent des classements
cohérents (Spearman 0,891). Un agent d'arbitrage sélectionne ensuite, sous un budget de 5\,M\$, 11 projets en
respectant le classement, les montants admissibles et un plafond sectoriel de 35\,\%. Les analyses de sensibilité
montrent que les décisions résistent aux paramètres de la méthode mais dépendent fortement de la simulation, ce qui
est la principale limite de la démonstration.

Le système est déployé dans une application Streamlit réutilisable : tout jeu de projets peut y être importé et
évalué avec ses propres critères, poids, filtres et budget. Un agent rédacteur et un conseiller conversationnel,
fondés sur Gemini, n'effectuent aucun calcul : ils appellent des fonctions Python déterministes et chaque nombre de
leurs réponses est vérifié automatiquement.

\medskip
\noindent\textbf{Mots-clés :} risque de crédit, probabilité de défaut, EBM, LightGBM, calibration, fuite de données,
aide multicritère à la décision, AHP, ELECTRE III, ELECTRE TRI, PROMETHEE II, sac à dos, agents LLM, Streamlit.
